\cotefile{21}
\scanpage{17}

% vertical rule at the left of the four lines
Si $X$ est loc noeth \unsure{quelct|quelcque|quelcq|quelconque} seulement, on dit que $F$ a la propriété $(P_1)$ (\unsure{resp\textsuperscript{t}|resp.} $(P_2)$). en un pt $x\in X$, si $\mathrm{Loc}_x(F)$ sur $\mathrm{Loc}_x(X)$ a la propriété $(P_1)$ resp. $(P_2)$.

\uline{Proposition} 3\quad Soit $X$ un espace loc. noeth. spectral, $Y$ une partie fermée de $X$, $F$ un faisceau d'ens sur $X$. Conditions équivalentes

(i) $F$ a la propriété $P_1$ (resp. $P_2$) rel à $Y$

(ii) Pour tt $x\in Y$, $F$ a la propriété $P_1$ (resp. $P_2$) en $x$.

\uline{Démonstration}\quad On peut supposer $X$ affine.

\deleted{\unsure{a)|$\alpha$)} \uline{\unsure{Cas de}} $(P_1)$} Prouvons d'abord le

\uline{Corollaire} 4\quad Si $F$ satisfait à la condition $P_1$ $(P_2)$ rel. à $Y$, alors pour tt ouvert $X'\subset X$ et toute partie fermée $Y'$ de $X'$ telle que $Y'\subset X'\cap Y$, $F':=F\mid X'$ satisfait $P_1$ $(P_2)$ rel à $Y'$.

% slanted
\margin{n'utilise pas l'hyp loc. noeth. sur $X$}

\inserted{Soit $U'=X-Y$.} On peut supposer évidemment $X'=X$. Soit $U$ un ouvert quelconque de $X$,

% archive sheet number: 16
